% Pista ana-25s-q3c · Anàlisi
% Procedència: PAU Catalunya 2025 · Convocatòria de setembre · Sèrie 3 · Exercici 3
\documentclass[11pt,a4paper]{article}
\usepackage{pista}
\pistainfo{Catalunya 2025}{Convocatòria de setembre}{Sèrie 3}

\begin{document}

\section*{\textcolor{blauPista}{Exercici 3}}

\noindent Una empresa de calçat esportiu ha creat una sabatilla amb amortiment per a minimitzar les lesions per sesamoïditis. Els beneficis generats per la venda d'aquest producte, en milers d'euros, segueixen una funció de la forma $f(x) = ax^{3} + bx^{2} + cx$, on $x$ són els anys transcorreguts des que la sabatilla va sortir a la venda i $a$, $b$ i $c$ són constants reals.

\subsection*{\textit{c)} \normalfont Calculeu els valors de $a$, $b$ i $c$ sabent que el primer any es van obtenir el màxim de beneficis, amb un valor de 8\,000 euros, i que el segon any va haver-hi un punt d'inflexió en els beneficis. \hfill\textnormal{[1 punt]}}

\pista{Tradueix cadascuna de les tres dades de l'enunciat a una condició matemàtica sobre $f$ o les seves derivades: ``el primer any hi ha el màxim de beneficis, amb valor 8'' (ves amb compte amb les unitats: són milers d'euros) i ``el segon any hi ha un punt d'inflexió''. Obtindràs un sistema de tres equacions amb tres incògnites ($a$, $b$, $c$).}

\end{document}
